% =============================================================================
% Appendix J: Sweep and Latency Measurement Protocols
% ORIUS/DC3S Battery-Safety Thesis
% =============================================================================
\chapter{Sweep and Latency Measurement Protocols}
\label{app:sweep-latency}

This appendix documents the exact protocols used for hyperparameter sensitivity sweeps, latency benchmarking, and fault severity sweeps in the ORIUS/DC3S battery-safety thesis. All procedures are reproducible via the referenced scripts.

\section{Hyperparameter Sweep Protocol}

\subsection{Grid Parameters}

The sensitivity sweep varies three DC3S hyperparameters over a full factorial grid:

\begin{itemize}
\item \texttt{alpha0}: Nominal miscoverage rate; values \texttt{0.05}, \texttt{0.10}, \texttt{0.15}.
\item \texttt{ph\_lambda}: Page-Hinkley drift detector sensitivity; values \texttt{3.0}, \texttt{5.0}, \texttt{7.0}.
\item \texttt{kappa\_drift\_penalty}: Drift penalty in $\kappa$-inflation; values \texttt{0.3}, \texttt{0.5}, \texttt{0.7}.
\end{itemize}

This yields $3 \times 3 \times 3 = 27$ grid points. Each point is run over scenarios \texttt{drift\_combo} and \texttt{dropout}, with seeds \texttt{11}, \texttt{22}, \texttt{33}, \texttt{44}, \texttt{55}, for a total of $27 \times 2 \times 5 = 270$ scenario-seed combinations per grid point.

\subsection{CLI Commands}

The sweep is executed via:

\begin{verbatim}
python scripts/run_sensitivity_sweeps.py \
  --out-dir reports/publication \
  --horizon 168 \
  --scenarios drift_combo,dropout \
  --seeds 11,22,33,44,55 \
  --alpha0 0.05,0.10,0.15 \
  --ph-lambda 3.0,5.0,7.0 \
  --kappa-drift-penalty 0.3,0.5,0.7 \
  --n-bootstrap 1000 \
  --confidence 0.95
\end{verbatim}

\subsection{Output Artifacts}

\begin{itemize}
\item \path{reports/publication/sensitivity_grid.csv} --- Full grid results (all rows).
\item \path{reports/publication/sensitivity_summary_ci.csv} --- Bootstrap confidence intervals by tuple $(\mathrm{scenario}, \mathrm{controller}, \alpha_0, \lambda_{\mathrm{PH}}, \kappa_{\mathrm{drift}})$.
\item \path{reports/publication/sweeps/configs/} --- Per-grid-point YAML configs; each filename records the sweep setting tuple.
\item \path{reports/publication/sweeps/runs/<tag>/} --- Per-run outputs such as \path{dc3s_main_table.csv}, \path{dc3s_fault_breakdown.csv}, \path{calibration_plot.png}, and \path{violation_vs_cost_curve.png}.
\end{itemize}

\subsection{Config Override Mechanism}

Each grid point sets \path{GRIDPULSE_DC3S_CONFIG} to the corresponding YAML
path. The base config \path{configs/dc3s.yaml} is loaded and overridden with
the grid values before each run.

\section{Latency Benchmark Protocol}

\subsection{Measurement Method}

The latency benchmark uses Python's \texttt{time.perf\_counter\_ns()} for high-resolution timing. Each component is measured as follows:

\begin{enumerate}
\item \textbf{Warmup}: Execute the callable \texttt{warmup} times (default 200) without recording. This reduces cold-cache effects.
\item \textbf{Measurement}: Execute the callable \texttt{iterations} times (default 10,000), recording elapsed nanoseconds per call.
\item \textbf{Statistics}: Compute mean, median, p95, p99, and max in milliseconds.
\end{enumerate}

\subsection{CLI Command}

\begin{verbatim}
python scripts/benchmark_dc3s_steps.py \
  --iterations 10000 \
  --warmup 200 \
  --out reports/dc3s_latency_benchmark.json
\end{verbatim}

\subsection{Benchmarked Components}

\begin{itemize}
\item \texttt{compute\_reliability} --- OQE reliability scoring.
\item \texttt{page\_hinkley\_update} --- Drift detector update.
\item \texttt{build\_uncertainty\_set} --- Linear inflation (RAC-Cert).
\item \texttt{build\_uncertainty\_set\_kappa} --- $\kappa$-inflation variant (if available).
\item \texttt{repair\_action} --- Shield projection.
\item \texttt{full\_step\_linear} --- Full DC3S step (linear inflation).
\item \texttt{full\_step\_kappa} --- Full DC3S step ($\kappa$ inflation).
\end{itemize}

\subsection{Output Artifacts}

\begin{itemize}
\item \texttt{reports/dc3s\_latency\_benchmark.json} --- Full JSON with per-component stats and environment metadata.
\item \texttt{reports/publication/dc3s\_latency\_summary.csv} --- CSV summary for tables (component, mean\_ms, median\_ms, p95\_ms, p99\_ms, max\_ms).
\end{itemize}

\section{Fault Severity Sweep}

The fault severity sweep varies dropout rate (or other fault intensity) to produce violation-rate and severity curves. It is integrated into the CPSBench runner via \texttt{include\_fault\_sweep=True}. For the main sensitivity sweep, \texttt{include\_fault\_sweep=False} is used to keep runtime manageable.

Fault sweep outputs include:
\begin{itemize}
\item \texttt{fig\_true\_soc\_violation\_vs\_dropout.png} --- Violation rate vs.\ dropout.
\item \texttt{fig\_true\_soc\_severity\_p95\_vs\_dropout.png} --- P95 severity vs.\ dropout.
\end{itemize}

These are generated by the CPSBench pipeline when fault sweep is enabled.
